\documentclass[11pt,letterpaper]{article}
\usepackage{amsmath,amssymb,amsthm}
\newtheorem{proposition}{Proposition}
\newtheorem{example}[proposition]{Example}
\newcommand{\N}{\mathbb{N}}
\newcommand{\F}{\mathcal{F}}

\title{Generating Functions for Counting}
\author{Rosa Delgado}
\date{April 5, 2023}

\begin{document}
\maketitle

\section{Ordinary generating functions}
To a sequence $(a_n)_{n\ge0}$ we associate the formal power series $A(x)=\sum_{n=0}^{\infty}a_nx^n$. Operations on sequences become operations on series: shifting, summing and convolving.

\subsection{The Fibonacci numbers}
Let $F_0=0$, $F_1=1$ and $F_{n+2}=F_{n+1}+F_n$. Multiplying the recurrence by $x^{n+2}$ and summing over $n\ge0$ gives $F(x)-x=xF(x)+x^2F(x)$, so
\begin{equation}\label{eq:fib}
  F(x)=\frac{x}{1-x-x^2}.
\end{equation}
Partial fractions yield Binet's formula
\[
  F_n=\frac{1}{\sqrt5}\Bigl[\Bigl(\frac{1+\sqrt5}{2}\Bigr)^{n}-\Bigl(\frac{1-\sqrt5}{2}\Bigr)^{n}\Bigr].
\]

\subsubsection{A consequence}
Since $|(1-\sqrt5)/2|<1$, the ratio $F_{n+1}/F_n$ converges to the golden ratio $\varphi=(1+\sqrt5)/2$.

\subsection{Binomial coefficients}
By the binomial theorem, $(1+x)^m=\sum_{k=0}^{m}\binom{m}{k}x^k$. Comparing coefficients in $(1+x)^m(1+x)^n=(1+x)^{m+n}$ gives Vandermonde's identity
\begin{equation}\label{eq:vdm}
  \sum_{k=0}^{r}\binom{m}{k}\binom{n}{r-k}=\binom{m+n}{r}.
\end{equation}

\begin{proposition}
The number of compositions of $n$ into parts from a set $S\subseteq\N$ has generating function $1/(1-\sum_{s\in S}x^s)$.
\end{proposition}

\begin{example}
Taking $S=\{1,2\}$ in the proposition gives $1/(1-x-x^2)$, so the number of compositions of $n$ into ones and twos is $F_{n+1}$, in line with \eqref{eq:fib}.
\end{example}

\section{Catalan numbers}
The Catalan numbers $C_n=\frac{1}{n+1}\binom{2n}{n}$ count, among many other things, balanced bracketings. They satisfy $C(x)=1+xC(x)^2$, and hence
\[
  C(x)=\frac{1-\sqrt{1-4x}}{2x}.
\]
The first values are $1,1,2,5,14,42,\dots$, and the sum $\sum_{i=0}^{n-1}C_iC_{n-1-i}=C_n$ is a special case of the product rule used to derive \eqref{eq:vdm}.

\end{document}
